\documentclass[10pt,letterpaper]{article}

% ---------- Paquetes ----------
\usepackage[margin=0.55in,top=0.5in,bottom=0.5in]{geometry}
\usepackage{titlesec}
\usepackage{enumitem}
\usepackage{hyperref}
\usepackage{fontspec}
\usepackage{xcolor}
\usepackage{tabularx}

% ---------- Fuente ----------
\setmainfont{Latin Modern Sans}

% ---------- Colores ----------
\definecolor{darkgray}{RGB}{40,40,40}
\definecolor{accent}{RGB}{20,60,110}

% ---------- Config de hipervínculos ----------
\hypersetup{
    colorlinks=true,
    urlcolor=accent,
    linkcolor=accent
}

% ---------- Estilo de títulos de sección ----------
\titleformat{\section}
  {\normalfont\large\bfseries\color{accent}}
  {}{0em}{}
  [\vspace{-4pt}{\color{accent}\hrule height 1pt}\vspace{1pt}]
\titlespacing{\section}{0pt}{6pt}{2pt}

% ---------- Sin numeración de página ----------
\pagenumbering{gobble}
\setlength{\parindent}{0pt}

% ---------- Listas compactas ----------
\setlist[itemize]{leftmargin=1.1em, itemsep=0pt, topsep=1pt, parsep=0pt}
\setlength{\emergencystretch}{2em}

\begin{document}

% ================= ENCABEZADO =================
\begin{center}
    {\Large \bfseries \color{darkgray} MARIO ALBERTO GUTIÉRREZ COVIÁN}\\[2pt]
    {Full Stack Software Engineer}\\[4pt]
    \small
    +52 (312) 103 1850 \;$\vert$\;
    \href{mailto:marioagtzc@hotmail.com}{marioagtzc@hotmail.com} \;$\vert$\;
    \href{https://github.com/MarioAGtzC}{GitHub} \;$\vert$\;
    \href{https://www.linkedin.com/in/marioagtz/}{LinkedIn} \;$\vert$\;
    \href{https://mario-gutierrez.com/}{Profile}
\end{center}

% ================= RESUMEN PROFESIONAL =================
\section*{Professional Summary}
Software development with experience in backend APIs and microservices using Node.js.
Experienced in designing and deploying serverless architectures on AWS, using services like
Lambda, Lex, API Gateway, DynamoDB, S3 and CloudWatch. Comfortable working with both SQL
and NoSQL databases.

% ================= LOGROS CLAVE =================
\section*{Key Achievements}
\begin{itemize}
    \item Implemented Amazon Lex conversational bots integrated with Connect and Lambda, automating customer interactions and reducing manual support workload.
    \item Enhanced IVR call flows and routing logic in Amazon Connect, improving user experience and reducing call handling time.
    \item Led backend improvements and bug fixes that increased system stability and reduced production issues.
    \item Collaborated on full-stack features using React and Redux, improving data flow efficiency and frontend performance.
    \item Built a Python/Django microservice that generated vector embeddings from unstructured files, enabling semantic search and improving search result relevance while reducing query response time through optimized data retrieval.
\end{itemize}

% ================= HABILIDADES TÉCNICAS =================
\section*{Technical Skills}
\begin{itemize}
    \item \textbf{Languages:} TypeScript (4 years), JavaScript (5 years), Python (2 years), Java (1 year)
    \item \textbf{Backend:} Node.js, Express.js, Nest.js, REST APIs, Django
    \item \textbf{Frontend:} React.js, Next.js, Redux, Tailwind CSS, HTML5, CSS3
    \item \textbf{Databases:} MySQL, MongoDB, PostgreSQL, DynamoDB, Redis, Firebase
    \item \textbf{Cloud \& Serverless:} AWS (Lambda, Lex, CloudWatch, API Gateway, Connect, S3), Docker
    \item \textbf{Other:} Git, Postman, Swagger, Agile/Scrum, Microservices Architecture, AI (GitHub Copilot, Claude, ChatGPT)
\end{itemize}

% ================= EXPERIENCIA PROFESIONAL =================
\section*{Professional Experience}

\textbf{Backend Software Engineer} \hfill \textit{April 2022 -- Current}\\
\textit{TechBridge IT | Remote}
\begin{itemize}
    \item Developed backend systems with TypeScript for robust, maintainable software.
    \item Designed and implemented Lex bots to enhance customer experience.
    \item Managed Lambda updates for better performance and scalability.
    \item Optimized call routing and workflows in Connect.
    \item Conducted performance analysis and testing on the Connect platform.
    \item Provided bilingual support and refined Spanish language options.
    \item Microservices maintenance in Python and Java.
    \item Managed an AI-powered context management microservice utilizing AWS Bedrock to automatically generate intelligent summaries of chat conversations between guests and agents.
\end{itemize}

\vspace{1pt}
\textbf{Full-Stack Developer} \hfill \textit{November 2020 -- April 2022}\\
\textit{Tata Consultancy Services | Remote}
\begin{itemize}
    \item Collaborated with front-end developers to design functional, cohesive code, enhancing user experience.
    \item Conducted data analysis to create efficient APIs for seamless system communication.
    \item Built new sections with React templates, optimizing data retrieval from AWS APIs using Redux.
    \item Set up development and production environments with AWS toolsets for streamlined deployment and scalability.
    \item Designed and managed database structures to ensure smooth information flow and data integrity.
\end{itemize}

\vspace{1pt}
\textbf{Full-Stack Web Developer} \hfill \textit{February 2020 -- November 2020}\\
\textit{Kioscomático | On-site, Colima}
\begin{itemize}
    \item Developed and maintained web applications with PHP and Node.js, focusing on performance, security, and scalability.
    \item Implemented RESTful APIs using Node.js, Express, and MySQL for backend development.
    \item Created a custom React components library to enhance code reusability and maintainability.
    \item Utilized React Hooks and Redux for consistent state management and user interface experience.
    \item Designed and developed mobile applications with Flutter for cross-platform compatibility.
    \item Integrated Bloc pattern for efficient state management in Flutter applications.
\end{itemize}

% ================= EDUCACIÓN =================
\section*{Education}
\textbf{Bachelor of Software Engineering} \hfill \textit{Graduated August 2020}\\
Universidad de Colima -- Facultad de Telemática

% ================= IDIOMAS =================
\section*{Languages}
\begin{itemize}
    \item \textbf{Spanish:} Native
    \item \textbf{English:} Professional Working Proficiency \href{https://cert.efset.org/en/SdWPH8}{(C1)}
\end{itemize}

\end{document}